Community-based practice is the accretion, translation, and distillation
of many voices across communities---who have shown me how to do this
work, what the work is, and how to be in community more gracefully and
more radically. For the last few years, I have been incredibly lucky to
have space to think, build, and scheme abundantly. The material traces
of this space can be found in the shared PhD workspace, SCI 211, where
projects in progress join projects past: crochet blankets, acoustic
dishes, spray paint, soldering irons, origami kites, a table full of
wood shavings and collage supplies, cables upon cables upon cables. This
space overflows in ways that barely hint at the joys and struggles of
creative research, scholarly rigor, personal and collective trials, and
the occasional faculty--PhD dance-off. Like SCI 211, an acknowledgments
page may appear to be a scattered collection, easily dismissed. But the
acknowledgments are the foundations of this text, the forces that make
possible all the efforts described within.

To my dissertation committee, who affirmed this work in all its forms
and pushed it forward with sharp, specific feedback, I am incredibly
honored to be part of your communities:

To Holly Willis, thank you for fourteen semesters as a steadfast mentor,
expertly moving from queer handmade praxis to generative AI to
bookbinding and back again. I'm grateful to be able to connect as
writers, makers, and researchers. You continually teach me how to look
and listen and offer with more nuance and grace. Thank you for expanding
my research practice and describing it with such clarity, even when I'd
call it a flailing tube man. Thank you for being a subtle mountain-mover
and fierce advocate for students and for weird art.

To Tara McPherson, thank you for showing me how to live my values as an
academic and citizen, how to create new tools for thinking differently,
and how to teach and think with both intellectual precision and absolute
whole-heartedness. Your advice is always spot-on and delivered with all
necessary dry wit. Thank you for modeling polymathic, polyvalent
feminism in the software lab, in the classroom, in scholarly journals,
and in the streets.

To Rita Raley, thank you for our in-depth discussions of publishing,
practice, and tactics---from technical details of natural language
processing to poetic affordances of electronic literature to forms of
resistance. I am grateful for the chance to think deeply with you across
all these modes and for your guidance in combining them.

I am also grateful to Scott Fisher and to Alice Gambrell for your
contributions and insights as part of my qualifying exams committee,
which laid the foundation that helped this work take shape and grow.
Thank you, Scott, for encouraging me to explore generative AI through
MEML projects and for making your classroom a space of broad
experimentation. Thank you, Alice, for inspiring discussions about
analog--digital publishing and for seeing something in my work that made
you encourage me to go to grad school all those years before I did.

It is hard to explain the magic of Code Collective, which has been both
the support system for this work and simultaneously its inspiration.
Thank you to everyone who has been a part of Code Collective, including
Katherine Yang, Elea Zhong, Ada Toydemir, Samir Ghosh, Garrett Flynn,
ender, Robby Feffer, Jay Borgwardt, Karen Abe, Amy Wu, Bita Tanavoli,
Katie Luo, Katie Liu, Johans Saldana Guadalupe, Nicole Carrara, Tim
Zhai, Tia Kemp, Kyle Ang, Adena Park, Kris Yuan, Remi Wedin, Dani
Takahashi, Ilona Bodnar, Emily Fuesler, Carey Crooks, Lee Thiboudeau,
Nico Pizzati, Ethan Kurzrock, Allison Ross, Bill Russell, Zachary Mann.
Whether or not your name is here, whether you joined as an in-person
regular or as a curious online visitor, thank you each for reminding me
how valuable Code Collective is as a community we make and keep making
together.

Thank you also to Dr.~Elizabeth Ramsey for your unflagging support of
that wild idea. You are the heart of MA+P and your office is a
sanctuary. I am grateful for your presence at every event, whether
supporting zine club, organizing the affective computing reading group,
or including me in your courses. Thank you for teaching me the secret to
fostering community (it's snacks) and for helping me foster it more.

Thank you to Dave Lopez for making all our iMAP ideas materialize, to
Sonia Seetharaman for making those ideas organized and beautiful, and to
Stacy Patterson for keeping us funded and well-fed. Thank you each for
navigating so much institutional bureaucracy on our behalf.

Abundant thanks go to the teachers who taught me about teaching,
mentorship, and community, among them Gabe Peters-Lazaro, Evan Hughes,
Safiya Noble, Jeanie Jo, and Anna Joy Springer. Thank you to Qianqian
Ye, who welcomed me into the p5.js community by letting me know I was
already part of it. Thank you to Brett Stalbaum, who knew I could code
before I did and pointed out a scrappy path. Thank you to Vicki Callahan
for generous dissertation support, both emotional and editorial, and a
zen approach to critical theory. Thank you to Steve Anderson for so much
heart-centered real-talk about academia. Thank you to Mark Marino,
co-author, collaborator, and friendtor, for introducing me with
unlimited bravado and including me in your brilliant schemes.

I also wish to acknowledge the additional years of PhD support I
received from USC Mellon Humanities in a Digital World Program, which
allowed me to expand this research and to design my own course, with
special thanks to Amy Braden, Will Cowan, and Sean Fraga.

At the Knowing Machines research project, I offer warmest thanks to Kate
Crawford for the opportunity to think with you across disciplines,
space-time dimensions, and formats. I am grateful for your generosity,
trust, encouragement, and discerning eye---as a mentor and editor, and
also as an artist and researcher. Thank you to Mike Ananny for rigorous,
kind, and considered perspectives on critical datasets and for your
generosity of spirit in convening MASTS and AIMS. Thank you to Vladan
Joler and Olivia Solis for thoughtful, big-hearted conversations and
insights into visual storytelling. Thank you to every member of Knowing
Machines for sharing a friendly, dynamic, interdisciplinary intellectual
community: Christo Buschek, Franny Corry, Melodi Dincer, Annie Dorsen,
Hannah Franklin, Ed Kang, Jake Karr, Sasha Luccioni, Will Orr, Jason
Schultz, Hamsini Sridharan, Jer Thorp, and Michael Weinberg.

I am so thankful for the generous support of the Akademie der Künste
Berlin. Thank you to Clara Herrmann and Nataša Vukajlovic for curatorial
genius, creative support, and friendship; to Maya Indira Ganesh and Nora
Khan for AI with care and without BS; to my fellow AI Anarchies Fellows
D'Andrade, Walla Capelobo, Sara Culmann, Petja Ivanova, Sonder, Tin
Wilke, Lara Fong Prosper, and especially Pedro Oliveira and Aarti Sunder
for countless process chats over coffee. I am grateful for the
hospitality of the staff and residents at Zentrum für Kunst und
Urbanistik, with particular thanks to Anita Rind, Mandy Seidler, Carmen
Santesmases, Vaas de Wit, Ariane Boucher, Noa Enghardt, and Agnes
Schyberg for all your tireless work turning a transitory public space
into a creative flourishing home.

At the Humboldt Institute for Internet and Society, thank you to Daniela
Dicks and Theresa Züger, to the AI \& Society Lab, and to everyone at
HIIG for welcoming me as a Fellow and for supporting the Intersectional
AI Toolkit in its development and hosting its launch. Thank you also to
everyone who has participated in Intersectional AI Toolkit zine-making
workshops from 2021--2024 at HIIG, USC, Mozilla Festival, re:publica,
transmediale Vorspiel (Halfsister), Akademie der Künste, ZK/U Berlin,
and Ars Electronica Founding Lab Summer School.

Thank you to the members of the USC Libraries' ``Visual Datasets for
Inclusive Research'' project, including Mike Jones, Bill Dotson, Curtis
Fletcher, Caroline Muglia, Hujefa Ali, and Manasa Rajesh for sharing
your knowledge, for encouraging intersectional perspectives, for
planting the seeds for future research projects, and for being the
happiest people on Zoom.

To Alena Giardina, thank you for making the roof over my head a
neighborly home and a launchpad for new endeavors. I am humbled and
continually inspired to pay it forward.

Where would I be without Fidelia Lam and Lisa Müller-Trede, my
weirdos-in-arms, and my fellow iMAP PhDs: Catherine Griffiths, Emilia
Yang, Triton Mobley, Szilvia Ruszev, Noa Kaplan, Romi Morrison, Ben
Nicholson, Kumi Iman, Luke Fischbeck, Sultan Sharrief, Dana Dal Bo,
Chantal Eyong, Faye Yan Zhang, Michelle Salinas, Selwa Sweidan, Curtis
Tam, EvB, Andrea Kim, Zeynep Abes, Ellie Schmidt, Sichong Xie, Halo
Starling, David de Rozas, David Kelley, Karl Baumann, Biayna Bogosian,
micha cárdenas, Brian Cantrell, Laura Cechanowicz, Behnaz Farahi, Todd
Furmanski, Aroussiak Gabrielian, Samantha Gorman, Juri Hwang, Geoff
Long, Clea Waite, Jeff Watson. Thank you for making SCI 211 a place to
belong in transdiciplinary eccentricity.

I am grateful to many creative, intellectual connections across and
beyond institutions, whose perspectives enhanced this work: Ashley
Dailey, Ariana Dongus, Katy Gero, Mario Guzman, JH, Shawné Michaelain
Holloway, Olivia Jack, Swantje Lichtenstein, Maurice Jones, Christine
Meinders, Miller Puckette, Lubna Rashid, Tiara Roxanne, Diana
Serbanescu, Lena Wegmann, and Xin Xin. For endless, uncategorizable,
above and beyond support, deepest thanks go to Åste Amundsen, cypress
masso, Emily Martinez, Leigh Montavon, Jason Perez, Susan Ring, Rob
Shafer, Tavia Stewart, and Isabel Wanger (all the information is on the
task). With profound love to my family, especially Anthony Ciston, Tony
Ciston, and Mary Jo Fisher.
